\begin{figure}[htbp]
\centering
\begin{tikzpicture}

% ========== Main Plot: Function value vs iteration ==========
\begin{axis}[
    name=mainplot,
    width=11cm, height=7cm,
    xlabel={Iteration $k$},
    ylabel={$f(\vecx_k)$},
    xmin=0, xmax=25,
    ymin=0, ymax=12,
    grid=major,
    grid style={UofSC50Black!30},
    title={\textbf{Gradient Descent Convergence}},
    title style={at={(0.5,1.02)}, font=\bfseries\large},
    legend pos=north east,
    legend style={font=\footnotesize, sharp corners, draw=UofSC70Black}
]

% Function: f(x) = 0.5*(x1^2 + 10*x2^2)
% Gradient: (x1, 10*x2)
% Starting from (3, 2): f_0 = 0.5*(9 + 40) = 24.5

% Large step size alpha = 0.15 (oscillatory convergence)
% For this function: x1_{k+1} = (1-0.15)*x1_k = 0.85*x1_k
%                    x2_{k+1} = (1-1.5)*x2_k = -0.5*x2_k (oscillates!)
\addplot[UofSCRose, very thick, mark=*, mark size=2pt, mark options={fill=UofSCRose}]
    coordinates {
        (0, 10.5) (1, 5.36) (2, 2.93) (3, 1.75) (4, 1.14)
        (5, 0.80) (6, 0.59) (7, 0.46) (8, 0.37) (9, 0.30)
        (10, 0.25) (11, 0.21) (12, 0.18) (13, 0.16) (14, 0.14)
        (15, 0.12) (16, 0.10) (17, 0.09) (18, 0.08) (19, 0.07)
        (20, 0.06) (21, 0.05) (22, 0.05) (23, 0.04) (24, 0.04)
    };
\addlegendentry{$\alpha = 0.15$ (large)}

% Medium step size alpha = 0.08 (good convergence)
\addplot[UofSCGarnet, very thick, mark=square*, mark size=2pt, mark options={fill=UofSCGarnet}]
    coordinates {
        (0, 10.5) (1, 7.14) (2, 4.91) (3, 3.41) (4, 2.40)
        (5, 1.71) (6, 1.23) (7, 0.90) (8, 0.66) (9, 0.49)
        (10, 0.37) (11, 0.28) (12, 0.21) (13, 0.16) (14, 0.13)
        (15, 0.10) (16, 0.08) (17, 0.06) (18, 0.05) (19, 0.04)
        (20, 0.03) (21, 0.02) (22, 0.02) (23, 0.02) (24, 0.01)
    };
\addlegendentry{$\alpha = 0.08$ (medium)}

% Small step size alpha = 0.02 (slow convergence)
\addplot[UofSCAtlantic, very thick, mark=triangle*, mark size=2pt, mark options={fill=UofSCAtlantic}]
    coordinates {
        (0, 10.5) (1, 9.52) (2, 8.64) (3, 7.85) (4, 7.14)
        (5, 6.51) (6, 5.94) (7, 5.42) (8, 4.96) (9, 4.54)
        (10, 4.16) (11, 3.82) (12, 3.51) (13, 3.23) (14, 2.98)
        (15, 2.75) (16, 2.54) (17, 2.35) (18, 2.18) (19, 2.02)
        (20, 1.88) (21, 1.75) (22, 1.63) (23, 1.52) (24, 1.42)
    };
\addlegendentry{$\alpha = 0.02$ (small)}

% Optimal value line
\draw[UofSCHorseshoe, thick, dashed] (axis cs:0,0) -- (axis cs:25,0);
\node[UofSCHorseshoe, font=\footnotesize, anchor=west] at (axis cs:25.2, 0) {$f^* = 0$};

\end{axis}

% ========== Inset: Zoomed view of convergence ==========
\begin{axis}[
    at={(mainplot.south east)},
    anchor=north east,
    xshift=-0.8cm,
    yshift=3.8cm,
    width=5cm, height=3.5cm,
    xlabel={$k$},
    ylabel={$f(\vecx_k)$},
    xlabel style={font=\footnotesize},
    ylabel style={font=\footnotesize},
    xmin=10, xmax=24,
    ymin=0, ymax=0.6,
    grid=major,
    grid style={UofSC50Black!30},
    tick label style={font=\tiny},
    title={\footnotesize\textbf{Zoomed: $k \geq 10$}},
    title style={at={(0.5,1.0)}, font=\footnotesize\bfseries}
]

% Large step size (zoomed)
\addplot[UofSCRose, thick, mark=*, mark size=1.5pt, mark options={fill=UofSCRose}]
    coordinates {
        (10, 0.25) (11, 0.21) (12, 0.18) (13, 0.16) (14, 0.14)
        (15, 0.12) (16, 0.10) (17, 0.09) (18, 0.08) (19, 0.07)
        (20, 0.06) (21, 0.05) (22, 0.05) (23, 0.04) (24, 0.04)
    };

% Medium step size (zoomed)
\addplot[UofSCGarnet, thick, mark=square*, mark size=1.5pt, mark options={fill=UofSCGarnet}]
    coordinates {
        (10, 0.37) (11, 0.28) (12, 0.21) (13, 0.16) (14, 0.13)
        (15, 0.10) (16, 0.08) (17, 0.06) (18, 0.05) (19, 0.04)
        (20, 0.03) (21, 0.02) (22, 0.02) (23, 0.02) (24, 0.01)
    };

% Small step size (zoomed)
\addplot[UofSCAtlantic, thick, mark=triangle*, mark size=1.5pt, mark options={fill=UofSCAtlantic}]
    coordinates {
        (10, 4.16) (11, 3.82) (12, 3.51) (13, 3.23) (14, 2.98)
        (15, 2.75) (16, 2.54) (17, 2.35) (18, 2.18) (19, 2.02)
        (20, 1.88) (21, 1.75) (22, 1.63) (23, 1.52) (24, 1.42)
    };

\end{axis}

% ========== Annotation: Update rule ==========
\node[font=\small, align=center] at (5.5,-0.8) {
    \textbf{Gradient Descent:} $\vecx_{k+1} = \vecx_k - \alpha \nabla f(\vecx_k)$
};

\node[font=\footnotesize, align=left, text width=6cm] at (9.5,-1.5) {
    $f(\vecx) = \frac{1}{2}(x_1^2 + 10x_2^2)$\\[2pt]
    Initial: $\vecx_0 = (3, 2)$, $f(\vecx_0) = 10.5$
};

\end{tikzpicture}
\caption{Convergence of gradient descent for a quadratic function $f(\vecx) = \frac{1}{2}(x_1^2 + 10x_2^2)$ with different step sizes $\alpha$. A large step size (Rose, $\alpha = 0.15$) converges fastest initially but may oscillate. A medium step size (Garnet, $\alpha = 0.08$) provides stable, efficient convergence. A small step size (Atlantic blue, $\alpha = 0.02$) guarantees convergence but is slow. The inset shows that after $k=10$ iterations, larger step sizes reach lower function values. The optimal step size balances speed and stability, typically $\alpha < 2/L$ where $L$ is the Lipschitz constant of the gradient.}
\label{fig:gradient_descent}
\end{figure}
